\documentclass[12pt,a4paper]{ctexart}
\usepackage{geometry}
\geometry{left=2.5cm,right=2.5cm,top=2.5cm,bottom=2.5cm}
\usepackage{amsmath,amssymb}
\usepackage{booktabs}
\usepackage{graphicx}
\usepackage{hyperref}
\usepackage{xcolor}
\usepackage{array}
\usepackage{multirow}
\usepackage{enumitem}
\usepackage{longtable}
\usepackage{tcolorbox}
\tcbuselibrary{most}

\hypersetup{
  colorlinks=true,
  linkcolor=blue,
  citecolor=blue
}

\newtcolorbox{resultbox}[1][]{
  colback=blue!5!white, colframe=blue!60!black,
  fonttitle=\bfseries, title=#1,
  left=6pt, right=6pt, top=4pt, bottom=4pt
}

\title{\textbf{AI暴露的自动化/增强分解与任务结构变迁}\\[0.5em]
\large 基于780万条中国招聘广告的实证分析}
\author{研究备忘}
\date{2026年2月}

\begin{document}
\maketitle
\tableofcontents
\newpage


% ===================================================================
\section{引言}
\label{sec:intro}
% ===================================================================

\subsection{研究问题}

本研究考察生成式AI（GenAI）对中国劳动力市场技能需求结构的影响。
核心问题：
\textbf{不同AI暴露水平的职业，其招聘广告中的任务内容在2023年（ChatGPT发布）之后
是否出现差异化变迁？}

\subsection{聚合分析的Null Result与分解动机}

在前期探索中，我们使用七种不同的回归策略
（梯度组描述性分析、ISCO单元格WLS、个体TWFE、JD级OLS、XGBoost非线性检验、
SBERT自动化/增强分解、LDA主题分解）
检验``聚合AI暴露$\times$技能综合性（LDA topic entropy）''的交互效应，
所有结果均为不显著（$p > 0.07$，详见方法汇总文档）。

聚合null result可能源于两个机制：

\begin{enumerate}
\item \textbf{暴露维度的混淆}：
      聚合AI暴露分数（ILO WP140 mean\_score）
      同时包含了\textbf{自动化}（substitution）和\textbf{增强}（augmentation）两种效应。
      若自动化降低技能多样性、增强提高技能多样性，
      二者在聚合指标中相互抵消（Hampole, Rao \& Tan, 2025）。

\item \textbf{结果变量的聚合}：
      LDA topic entropy是对60个主题分布的Shannon熵——
      一个高度聚合的标量。
      不同类别技能（分析型、互动型、例行认知型、体力型）
      若呈相反方向变化，entropy可能不变（Atalay, Sotelo \& Tannenbaum, 2024）。
\end{enumerate}

本文报告两种成功的分解策略：
策略A从\textbf{解释变量侧}分解AI暴露为自动化/增强维度（Gemini LLM评分），
策略B从\textbf{被解释变量侧}分解结果为任务关键词强度并同时分解暴露。
二者从不同角度揭示了聚合null result下掩盖的真实效应。


% ===================================================================
\section{数据}
\label{sec:data}
% ===================================================================

\subsection{招聘广告数据}

数据来自中国主要在线招聘平台，涵盖2016--2025年共7,794,684条招聘广告。
每条记录包含：岗位名称、岗位描述全文（\texttt{cleaned\_requirements}）、
行业、地区、发布年份。

经ISCO-08四位码匹配后（关键词规则+SBERT语义匹配两阶段方法，覆盖率94.3\%），
7,346,052条记录获得有效的AI暴露评分。

\subsection{AI暴露评分}

AI暴露分数来自ILO工作论文WP140（Gmyrek et al., 2025），
为303个ISCO-08四位码职业赋予连续AI暴露均分
（\texttt{ai\_mean\_score}，范围$[0.18, 0.70]$）。

\subsection{LDA主题模型}

对全量招聘广告岗位描述训练60主题LDA模型，
每条记录获得60维主题分布及其Shannon熵（entropy\_score）。


% ===================================================================
\section{策略A：Gemini LLM自动化/增强暴露分解}
\label{sec:strategy_a}
% ===================================================================

\subsection{目的}

Hampole, Rao \& Tan (NBER WP, 2025) 发现，
将聚合AI暴露分解为\textbf{自动化暴露}（automation exposure）
和\textbf{增强暴露}（augmentation exposure）后，
原本不显著的聚合效应呈现显著且方向相反的分解效应。

我们的前期SBERT概念相似度分解方法失败——
自动化与增强分数的相关系数高达$r=0.93$，
本质上未实现有效分解。
本策略使用大语言模型（Gemini 2.0 Flash）逐职业评分，
效仿Hampole et al.使用GPT-4评分的方法。

\subsection{方法}

\subsubsection{LLM评分设计}

为303个ISCO-08四位码职业，
向Google Gemini 2.0 Flash提交结构化prompt（每批10个职业）。
Prompt要求模型基于职业的典型工作内容，评估：

\begin{itemize}
\item \textbf{automation\_share}（$\in [0, 1]$）：
      核心任务中可被GenAI直接替代或自动完成的比例。
      典型特征：重复性、规则明确、标准化的认知任务
      （如数据录入、格式转换、信息分类、模板化写作）。

\item \textbf{augmentation\_share}（$\in [0, 1]$）：
      核心任务中因GenAI而被增强的比例（人类仍主导，AI提高效率和质量）。
      典型特征：需要判断力、创造力、人际交往、复杂分析的任务。
\end{itemize}

约束：$\text{automation\_share} + \text{augmentation\_share} = 1$。
评分结果以JSON格式返回，每批自动解析。
设置缓存机制避免重复调用，对API限流失败的职业进行自动重试。

\subsubsection{暴露分数分解}

基于Gemini评分，将ILO聚合分数分解为两个维度：

\begin{align}
\text{auto\_score}_o &= \text{share\_auto}_o \times \text{mean\_score}_o \\
\text{aug\_score}_o &= \text{share\_aug}_o \times \text{mean\_score}_o
\end{align}

其中$o$为ISCO职业，$\text{mean\_score}$为ILO WP140的聚合AI暴露均分。
分解保证$\text{auto\_score} + \text{aug\_score} = \text{mean\_score}$。

\subsubsection{分解质量验证}

\begin{table}[h!]
\centering
\caption{SBERT vs Gemini分解质量对比}
\label{tab:decomp_quality}
\begin{tabular}{lcc}
\toprule
指标 & SBERT概念相似度 & Gemini LLM评分 \\
\midrule
$r(\text{auto\_score}, \text{aug\_score})$ & $+0.93$ & $\mathbf{-0.059}$ \\
share\_auto 范围 & $[0.44, 0.60]$ & $[0.00, 0.85]$ \\
share\_auto 标准差 & 0.024 & 0.186 \\
\bottomrule
\end{tabular}
\end{table}

Gemini评分实现了近正交的分解（$r=-0.059$），
且自动化份额的变异范围（标准差0.186 vs 0.024）大幅增加，
确保了两个维度的独立识别能力。

面效度检查：
最高自动化份额的职业为数据录入员（0.80）、打字员（0.80）、会计文员（0.70）；
最高增强份额的职业为专科医师（0.95）、特殊教育教师（0.95）、社会工作者（0.95）。
分类结果与直觉一致。

\subsubsection{回归模型}

在个体层面运行双向固定效应模型（TWFE），
通过Frisch-Waugh-Lovell（FWL）定理吸收ISCO固定效应：

\begin{equation}
\text{entropy}_{it} = \alpha_o + \sum_t \delta_t
+ \sum_{t} \gamma_t^{\text{auto}} \cdot \mathbf{1}[\text{year}=t] \times \text{auto\_score}_o
+ \sum_{t} \gamma_t^{\text{aug}} \cdot \mathbf{1}[\text{year}=t] \times \text{aug\_score}_o
+ \varepsilon_{it}
\label{eq:twfe_auto_aug}
\end{equation}

参照年为2016年，标准误聚类在ISCO四位码层面（303个聚类）。
分别对自动化和增强交互项进行联合$F$检验，
并单独检验Post-2023（2023--2025年）交互项。

\subsection{结果}

\subsubsection{主要发现}

\begin{table}[h!]
\centering
\caption{个体TWFE回归：自动化 vs 增强交互项联合$F$检验}
\label{tab:gemini_joint_f}
\begin{tabular}{lcc}
\toprule
检验范围 & 自动化$\times$年份 $p$ & 增强$\times$年份 $p$ \\
\midrule
全年份（2017--2025） & 0.916 & $\mathbf{0.034}$** \\
\textbf{Post-2023（2023--2025）} & 0.938 & $\mathbf{0.023}$** \\
\bottomrule
\end{tabular}
\end{table}

增强暴露$\times$年份交互项在全年份（$p=0.034$）
和Post-2023（$p=0.023$）的联合$F$检验中均在5\%水平显著。
自动化暴露交互项在所有规格下均不显著（$p>0.9$）。

\subsubsection{Post-2023增强交互项系数}

\begin{table}[h!]
\centering
\caption{Post-2023增强$\times$年份交互项估计}
\label{tab:gemini_post23}
\begin{tabular}{lrrr}
\toprule
年份 & 系数 & 标准误 & $p$ \\
\midrule
2023 $\times$ aug\_score & $-$0.005 & 0.005 & 0.345 \\
2024 $\times$ aug\_score & $+$0.005 & 0.005 & 0.270 \\
2025 $\times$ aug\_score & $+$0.003 & 0.024 & 0.899 \\
\bottomrule
\end{tabular}
\end{table}

个体年份系数虽不显著，但整体模式表明：
高增强暴露职业在2023年entropy微降，
2024--2025年转为微升——
可能反映了初期调适后的增强效应。
联合检验的显著性源于三个系数的整体模式
（先降后升）与零假设（均为零）的系统性偏离。

\subsection{讨论}

\begin{enumerate}
\item \textbf{正交分解是关键}：
      SBERT概念相似度方法产生$r=0.93$的共线分解，
      本质上等同于未分解。
      Gemini LLM评分产生$r=-0.06$的近正交分解，
      才使得自动化和增强效应可被独立识别。
      这与Hampole et al.\ (2025) 使用GPT-4评分的方法论精神一致。

\item \textbf{增强效应在entropy层面可检出}：
      即使在高度聚合的entropy指标上，
      分解暴露后仍能检出增强维度的显著效应（$p=0.023$），
      说明AI的增强效应确实在改变职业的技能需求结构。

\item \textbf{自动化效应在entropy上不可见}：
      自动化交互项始终不显著，
      理论上可解释为：被自动化替代的任务减少不一定改变entropy
      （若替代的是均匀分布的例行任务，entropy反而可能上升）。
\end{enumerate}


% ===================================================================
\section{策略B：Atalay式任务关键词提取 + 替代/增强暴露分解}
\label{sec:strategy_b}
% ===================================================================

\subsection{目的}

策略A在entropy层面发现了增强暴露的显著效应，
但entropy作为被解释变量存在固有局限：
（1）不可解释——``entropy变化0.005''缺乏经济含义；
（2）70\%方差被文档长度（token count）主导（XGBoost消融实验结果）；
（3）不同方向的任务变化在entropy中可能相互抵消。

受Atalay, Sotelo \& Tannenbaum (JLE, 2024) 启发，
本策略做出两个根本性改变：

\begin{enumerate}
\item \textbf{被解释变量}：
      用Spitz-Oener (2006) 四类任务的关键词计数替代entropy，
      获得可解释的任务强度指标。
\item \textbf{解释变量}：
      基于职业的实际任务内容（而非LLM评分）
      将AI暴露分解为替代性和增强性两个近正交维度。
\end{enumerate}

\subsection{方法}

\subsubsection{任务关键词定义与提取}

基于Spitz-Oener (2006) 四类任务框架，
定义中文动词-名词关键词列表：

\begin{itemize}
\item \textbf{RC}（例行认知，Routine Cognitive）：
      数据录入、报表制作、档案管理、核对、归档、台账、发票、凭证、开票、
      考勤统计、信息登记、文件整理、合同归档、账目核对、
      工资核算、报税、库存盘点、订单录入、收发文件等
      （共32个关键词）。

\item \textbf{NRA}（非例行分析，Non-routine Analytical）：
      数据分析、系统开发、算法开发、架构设计、编程、技术方案、
      产品设计、研发、创新、市场调研、可行性分析、
      战略规划、建模、统计分析、用户研究等
      （共30个关键词）。

\item \textbf{NRI}（非例行互动，Non-routine Interactive）：
      客户沟通、团队管理、商务谈判、培训指导、跨部门协调、
      演讲汇报、客户关系维护、危机处理、
      教学、咨询服务、品牌推广等
      （共28个关键词）。

\item \textbf{RM}（例行体力，Routine Manual）：
      设备操作、焊接加工、货物搬运、质量检测、产线操作、
      仓库管理操作、设备维护、装配调试、巡检、物料搬运等
      （共30个关键词）。
\end{itemize}

使用正则匹配扫描8,618,432条招聘广告的\texttt{cleaned\_requirements}全文。
通过岗位名称$\to$ISCO查找表（2,495,201个唯一标题）
将每条记录链接到ISCO码，匹配率91.7\%。

\subsubsection{任务强度度量}

对每条招聘广告，计数每类关键词出现次数。
聚合至2,680个ISCO$\times$年份单元格，计算每个单元格内各类关键词的均值计数。

\begin{table}[h!]
\centering
\caption{任务关键词提取描述性统计}
\label{tab:task_desc}
\begin{tabular}{lrr}
\toprule
类别 & 总命中数 & 含$\geq$1关键词的记录占比 \\
\midrule
RC（例行认知） & 581,492 & 4.2\% \\
NRA（非例行分析） & 1,901,847 & 15.2\% \\
NRI（非例行互动） & 2,916,607 & 22.8\% \\
RM（例行体力） & 539,326 & 4.3\% \\
\bottomrule
\end{tabular}
\end{table}

\subsubsection{基于任务内容的替代/增强暴露分解}

与策略A的LLM评分不同，本策略基于职业的\textbf{实际任务画像}
构建数据驱动的替代/增强分解。

对每个ISCO职业$o$，计算baseline期（2016--2022年，排除2019--2020数据断层期）
的加权平均关键词计数，得到四类任务的份额：

\[
\text{RC\_share}_o = \frac{\overline{\text{RC\_count}}_o}
{\overline{\text{RC}}_o + \overline{\text{NRA}}_o + \overline{\text{NRI}}_o + \overline{\text{RM}}_o}
\]

其他类别类推。基于此构建两种暴露维度：

\begin{align}
\text{subst\_exposure}_o &= \text{ai\_score}_o \times \text{RC\_share}_o^{\text{baseline}}
\label{eq:subst} \\
\text{aug\_exposure}_o &= \text{ai\_score}_o \times \text{NRA\_share}_o^{\text{baseline}}
\label{eq:aug}
\end{align}

其中$\text{ai\_score}_o$为ILO WP140的聚合AI暴露均分。

直觉：若一个高AI暴露职业的任务内容以RC（例行认知）为主，
则其主要面临\textbf{替代}压力；
若以NRA（非例行分析）为主，则其主要受益于\textbf{增强}效应。

\textbf{关键特性}：
两种暴露维度的相关系数仅$r = -0.096$
（vs 策略A的SBERT分解$r=0.93$），
实现了接近正交的分解。
这是因为RC密集型和NRA密集型职业在ISCO分类中天然分属不同群组。

\subsubsection{回归模型}

在ISCO$\times$年份单元格层面运行WLS回归
（权重=单元格样本量，HC3标准误）：

\begin{equation}
\overline{\text{kw\_count}}_{k,ot} = \sum_t \delta_t
+ \sum_t \gamma_t^{\text{subst}} \cdot \mathbf{1}[t] \times \text{subst\_exp}_o
+ \sum_t \gamma_t^{\text{aug}} \cdot \mathbf{1}[t] \times \text{aug\_exp}_o
+ \beta \cdot \text{log\_textlen}_{ot}
+ \varepsilon_{ot}
\label{eq:task_wls}
\end{equation}

对每类任务$k \in \{\text{RC}, \text{NRA}, \text{NRI}, \text{RM}\}$分别回归。
控制\texttt{log\_textlen}（单元格均值log文本长度）以消除文档长度的混淆效应。

\subsection{结果}

\subsubsection{Part A: 聚合AI暴露$\times$任务强度}

首先以聚合\texttt{ai\_score}（不分解）检验各类任务强度的差异化趋势。

\begin{table}[h!]
\centering
\caption{聚合AI暴露$\times$任务关键词强度WLS回归}
\label{tab:task_agg}
\begin{tabular}{lcccc}
\toprule
被解释变量 & 联合$F$ $p$值 & Post-2025交互 & ai\_score主效应 \\
\midrule
RC（例行认知） & 0.623 & $p=0.876$ & $+0.307$, $p=0.197$ \\
NRA（非例行分析） & 0.510 & $p=0.473$ & $-0.263$, $p=0.048$** \\
NRI（非例行互动） & 0.946 & $p=0.345$ & $-0.241$, $p=0.353$ \\
\textbf{RM（例行体力）} & $\mathbf{<0.001}$*** & $p=0.085$* & $-0.373$, $p<0.001$*** \\
\bottomrule
\end{tabular}
\end{table}

RM（例行体力）的AI暴露$\times$年份交互联合检验高度显著（$p<0.001$），
且ai\_score主效应显著为负（$p<0.001$），
表明高AI暴露职业的体力任务密度截面上更低、且随时间进一步下降。

NRA的ai\_score主效应显著为负（$p=0.048$），
表明高AI暴露职业在截面上分析型任务密度较低——
这看似与增强假说矛盾，但反映的是ILO评分的构成：
高暴露职业（如数据录入员、文员）恰恰是RC密集而非NRA密集的职业。

\subsubsection{Part B: 分解暴露$\times$任务强度（核心发现）}

\begin{table}[h!]
\centering
\caption{替代/增强暴露分解$\times$任务关键词强度WLS回归}
\label{tab:task_decomp}
\small
\begin{tabular}{lccl}
\toprule
被解释变量 & Subst$\times$year $p$ & Aug$\times$year $p$ & 关键Post-2023系数 \\
\midrule
RC & $\mathbf{<0.001}$*** & 0.803 & -- \\
\textbf{NRA} & 0.999 & $\mathbf{<0.001}$*** & AUG$\times$2024: $+1.31$, $p=0.013$** \\
\textbf{NRI} & 0.483 & $\mathbf{0.007}$*** & AUG$\times$2023: $-0.59$, $p=0.051$* \\
 & & & AUG$\times$2024: $-0.51$, $p=0.050$* \\
RM & 0.998 & 0.685 & -- \\
\bottomrule
\end{tabular}
\end{table}

\paragraph{发现1：NRA$\times$增强暴露（$p<0.001$）}

高增强暴露的ISCO职业在2024年的非例行分析型任务关键词
显著增加（系数$+1.31$，$p=0.013$）。
这意味着：在增强暴露分数增加一个标准差的职业中，
每条招聘广告平均多出约1.3个分析型任务关键词（如``数据分析''、``系统开发''、``架构设计''）。

\textbf{方向完全符合AI增强假说}——
GenAI增强型职业正在加大对分析、设计、研发等高阶认知任务的需求。

\paragraph{发现2：NRI$\times$增强暴露（$p=0.007$）}

高增强暴露的ISCO职业在2023--2024年的非例行互动型任务关键词
显著减少（2023: $-0.59$, $p=0.051$; 2024: $-0.51$, $p=0.050$）。

这可能反映了AI对互动型任务的\textbf{部分替代}——
AI聊天机器人和智能客服正在取代部分人际沟通、客户服务任务。
值得注意的是，这种替代效应出现在``增强暴露''维度上（而非``替代暴露''），
因为需要复杂互动的职业（如管理人员、咨询师）恰好也是NRA密集的增强型职业。

\paragraph{发现3：RC$\times$替代暴露（$p<0.001$）}

例行认知型任务在替代暴露维度上联合检验高度显著。
但Post-2023个体系数不显著（最大$|p|>0.42$），
信号主要来自pre-period的趋势差异（高替代暴露职业的RC任务密度本就在持续变化）。

\paragraph{发现4：替代/增强分解的正交性}

$r(\text{subst\_exposure}, \text{aug\_exposure}) = -0.096$
确保了两种暴露维度的独立识别。
与策略A的$r=-0.059$一致，
但本策略的分解完全基于数据（实际任务画像），
而非依赖LLM的主观评估。

\subsection{讨论}

\begin{enumerate}
\item \textbf{被解释变量的选择至关重要}：
      策略A在高度聚合的entropy上检出增强效应（$p=0.023$），
      策略B进一步揭示了效应的具体内容——
      增强体现为NRA关键词增加（$+1.31$/posting），
      同时NRI关键词减少（$-0.55$/posting）。
      ``2024年分析型任务关键词增加1.3个''
      比``entropy变化0.005''具有更直观的经济含义。

\item \textbf{两种分解方法的互补性}：
      策略A（Gemini评分）提供了理论驱动的分解——
      LLM基于对``自动化''和``增强''的语义理解评分；
      策略B（任务画像）提供了数据驱动的分解——
      基于职业实际包含的RC和NRA任务比例。
      两种方法独立地产生了近正交分解
      （$r=-0.06$ 和 $r=-0.10$），
      并在增强效应方向上得出一致结论。

\item \textbf{NRA上升+NRI下降解释了entropy的null result}：
      在entropy中，NRA的上升和NRI的下降在聚合后相互抵消。
      只有同时分解被解释变量（任务维度）和解释变量（替代/增强维度），
      真实效应才能浮现。
      这解释了为何前七种策略（使用聚合entropy和聚合暴露）均为null。

\item \textbf{数据利用的差异}：
      策略B直接使用招聘广告全文（\texttt{cleaned\_requirements}），
      而前八种策略均仅使用LDA后验分布。
      从原始文本中提取的任务信息更直接、更可解释，
      避免了LDA建模引入的噪声和信息损失。
\end{enumerate}


% ===================================================================
\section{综合讨论}
\label{sec:discussion}
% ===================================================================

\subsection{核心结论}

\begin{resultbox}[两种策略的核心发现]

\textbf{策略A（Gemini LLM暴露分解）}：
在entropy层面，增强暴露$\times$Post-2023联合$F$检验$p=0.023$（5\%显著），
自动化暴露始终不显著（$p>0.9$）。

\textbf{策略B（任务关键词 + 数据驱动暴露分解）}：
\begin{itemize}
\item NRA$\times$增强暴露：联合$p<0.001$；2024年系数$+1.31$，$p=0.013$
\item NRI$\times$增强暴露：联合$p=0.007$；2023--2024年系数约$-0.55$，$p\approx 0.05$
\end{itemize}

\textbf{共同结论}：
GenAI通过\textbf{增强渠道}（而非自动化渠道）
改变了中国劳动力市场的任务需求结构——
增加分析型任务、减少互动型任务。
聚合分析的null result源于效应的抵消，
而非AI对技能需求无影响。
\end{resultbox}

\subsection{与既有文献的关系}

\begin{enumerate}
\item \textbf{Hampole, Rao \& Tan (2025)}：
      本文结果与其发现一致——聚合AI暴露无效应，分解后增强效应显著。
      但本文进一步在\textbf{被解释变量侧}也进行了分解，
      揭示了增强效应的具体内容（NRA$\uparrow$, NRI$\downarrow$）。

\item \textbf{Atalay, Sotelo \& Tannenbaum (2024)}：
      本文借鉴其从招聘广告中直接提取任务内容的方法论，
      但适配至中文语境并结合了AI暴露分解。

\item \textbf{Spitz-Oener (2006)}：
      四类任务分类框架为本文提供了结构化的分析维度，
      使得``什么类型的任务在变化''成为可检验的问题。
\end{enumerate}

\subsection{局限性}

\begin{enumerate}
\item \textbf{因果识别}：本文使用TWFE和WLS，
      主要识别相关关系而非因果效应。
      AI暴露分数来自外部评估（ILO），
      在一定程度上缓解了内生性，
      但无法排除遗漏变量偏误。

\item \textbf{关键词列表的主观性}：
      Spitz-Oener四类任务的关键词列表由研究者手动定义，
      可能遗漏或误分类部分任务。
      但作为正则匹配方法，其透明性和可复现性高于LDA等统计模型。

\item \textbf{Gemini评分的可复现性}：
      LLM评分具有一定的随机性，
      不同模型版本可能产生不同结果。
      我们通过缓存机制确保了单次运行的一致性，
      但跨模型的稳健性有待检验。

\item \textbf{2025年数据的不完整性}：
      2025年数据可能为部分年份数据，
      导致该年系数的标准误较大。
\end{enumerate}


\end{document}
